\documentclass[11pt,a4paper]{article}

\usepackage[margin=2.8cm]{geometry}
\usepackage{amsmath,amssymb,amsthm,mathtools}
\usepackage{microtype}
\usepackage[hidelinks]{hyperref}
\usepackage[capitalize,noabbrev]{cleveref}

\newtheorem{theorem}{Theorem}
\newtheorem{lemma}[theorem]{Lemma}
\newtheorem{proposition}[theorem]{Proposition}
\newtheorem{corollary}[theorem]{Corollary}
\newtheorem{claim}[theorem]{Claim}
\theoremstyle{definition}
\newtheorem{remark}[theorem]{Remark}

\newcommand{\cFPS}{c_{\mathrm{FPS}}}
\newcommand{\NN}{\mathbb{N}}
\newcommand{\QQ}{\mathbb{Q}}
\newcommand{\RR}{\mathbb{R}}
\newcommand{\Eu}{\mathbb{E}}
\newcommand{\chip}{\chi'}
\DeclareMathOperator{\crN}{cr}

\title{A tighter constant in the chromatic-index lemma of\\
       Fox--Pach--Suk towards Albertson's conjecture}
\author{(draft) Marc Lelarge\thanks{\texttt{marc.lelarge@gmail.com}.}}
\date{\today}

\begin{document}

\begin{center}
{\Large\bfseries RETRACTED (2026-05-16)}\\[1em]
\fbox{\parbox{0.9\textwidth}{\itshape
The main theorem of this note is \textbf{false}.  The proof of
Lemma~6 (\emph{Case 2b}) silently assumed that the FPS monotonicity
of $f_{2b}(\eta,\delta)$ in $\eta$ on $[1/2, \eta_b(\delta))$ holds at
every $\delta > 1$.  In fact it holds only for $\delta \ge -3 + \sqrt{17}
\approx 1.12311$; at the proposed
$\delta_1 \approx 1.114907$, $f_{2b}^{\max}(\delta_1)
\approx 0.5654 > 9/16$, so $F(\delta_1) > F(9/8) = 9/16$ and no
improvement is obtained.\\[0.6em]
The corrected analysis shows that FPS's choice $\delta = 9/8$ is
structurally forced and that $F(\delta)$ attains its minimum at this
value, with $F(9/8) = 9/16$.  See the companion note
\texttt{sharpness\_9\_8.tex} for the corrected statement.\\[0.6em]
The original (false) draft is preserved below for historical record.
}}
\end{center}
\vspace{2em}

\maketitle

\begin{abstract}
Fox, Pach and Suk (SoCG~2025) recently proved Albertson's conjecture
$\crN(G)\ge\crN(K_k)$ for every $k$-chromatic graph $G$ on at most
$1.4(k-1)$ vertices. The central technical input is their Lemma~2.3,
which bounds the chromatic index of an auxiliary multigraph by
$(9/16+o(1))\,k$. The constant $9/16$ in turn arises from a three-case
optimisation (their Claim~3.7) parameterised by a degree threshold
$d:=\delta k$, set to $\delta=9/8$. We observe that this choice of
$\delta$ is not the analytic optimum of Claim~3.7's
case-bound function. Reoptimising over $\delta$ replaces $9/16$ by
\[
   F^{\star} \;=\; -\tfrac{1}{2}+\tfrac{2}{\sqrt3}\,
                   \cos\!\left(\tfrac{1}{3}\arccos\tfrac{3\sqrt3}{16}\right)
            \;\approx\; 0.557454,
\]
the unique real root in $(\tfrac12,\tfrac58)$ of $4x^3+6x^2-x-2=0$.
Every other ingredient of Fox--Pach--Suk's argument is checked to be
robust to this change. The resulting tightening of Lemma~2.3 is modest
in absolute terms ($\approx0.9\%$) but follows immediately from their
proof.
\end{abstract}

\section{Introduction}

Albertson conjectured~\cite{Albertson} that for every $k\ge1$ and every
graph $G$ with chromatic number $\chi(G)\ge k$,
\begin{equation}\label{eq:albertson}
  \crN(G) \;\ge\; \crN(K_k),
\end{equation}
where $\crN(\cdot)$ denotes the crossing number. The conjecture is
known for $k\le24$, with substantial recent progress: Albertson, Cranston
and Fox~\cite{ACF} proved the case $k\le12$; Bar\'at and
T\'oth~\cite{BaratToth} extended it to $k\le16$; Ackerman~\cite{Ackerman}
to $k\le18$; and Cranston~\cite{Cranston} to $k\le24$, with strong
restrictions on the possible counterexamples at $k\in\{25,26\}$.

Fox, Pach and Suk~\cite{FPS} prove~\eqref{eq:albertson} for every
$k$-chromatic graph on at most $1.4(k-1)$ vertices (and asymptotically
on $(1.64-o(1))\,k$ vertices) by a clever weak-immersion argument. The
crucial technical step is the following bound on the chromatic index of
an auxiliary multigraph $H$ constructed from $G$ via a semi-random
selection procedure (see~\Cref{sec:reconstruction} for the precise
setup):
\begin{equation}\label{eq:fpsbound}
  \chip(H_i) \;\le\; \bigl(\cFPS+o(1)\bigr)\,k_i,
  \qquad \cFPS=\tfrac{9}{16}.
\end{equation}
This is their Lemma~2.3. They derive the constant $\cFPS=9/16$ from a
three-case optimisation parameterised by a single free real
\emph{degree threshold} $\delta$; throughout~\cite{FPS} they fix
$\delta=9/8$ from the start.

In this note we observe that the choice $\delta=9/8$ in~\cite{FPS} is
suboptimal for the case bound. Re-optimising the three-case bound in
their Claim~3.7 with $\delta$ as a free parameter yields a smaller
constant.

\begin{theorem}\label{thm:main}
Let $F^{\star}$ be the unique real root in $\bigl(\tfrac12,\tfrac58\bigr)$
of $4x^3+6x^2-x-2=0$. Then~\eqref{eq:fpsbound} holds with
$\cFPS$ replaced by $F^{\star}$.

Equivalently, $\cFPS$ can be replaced by
\[
  F^{\star}
  \;=\; -\frac12 + \frac{2}{\sqrt3}\,
        \cos\!\left(\frac13\arccos\frac{3\sqrt3}{16}\right)
  \;=\; 0.5574537707\ldots
\]
\end{theorem}

\noindent
Since $F^{\star}<\tfrac{9}{16}=0.5625$, this is a strict improvement over
\cite[Lem.~2.3]{FPS}. The size of the improvement is
$\tfrac{9}{16}-F^{\star}\approx 5.05\cdot10^{-3}$, or about
$0.9\%$. The optimal degree threshold realising $F^\star$ is
\[
  \delta_1
  \;=\; -1+\frac{4}{\sqrt3}\,
        \cos\!\left(\frac13\arccos\frac{3\sqrt3}{16}\right)
  \;=\; 1.1149075414767\ldots,
\]
the unique real root in $(1,5/4)$ of the irreducible cubic
\begin{equation}\label{eq:cubic}
   \delta^3+3\delta^2-\delta-4 = 0.
\end{equation}

The improvement is intentionally small in scope: we change exactly one
numerical constant inside Fox--Pach--Suk's proof of their Lemma~2.3 and
verify (\Cref{sec:robustness}) that none of the other ingredients of
the proof is affected. Whether $F^{\star}$ itself is sharp -- or
whether a more substantive improvement of~\eqref{eq:fpsbound} is
possible by attacking the underlying Vizing--Gupta input -- remains
open (\Cref{rem:sharp}).

The downstream impact on the vertex-bound coefficients of
\cite[Thm.~1.2]{FPS} is straightforward but bookkeeping-heavy; we
isolate it as~\Cref{rem:downstream}.

\section{Reconstruction of \texorpdfstring{\cite{FPS}}{[FPS]}, Claim 3.7}
\label{sec:reconstruction}

We follow the notation of~\cite{FPS} throughout. Let $G$ be a
$k$-chromatic graph and write $\{V_1,\dots,V_r\}$ for the parts of the
Gallai decomposition of $G$, with $G[V_i]$ being $k_i$-critical and
$n_i:=|V_i|\ge 2k_i-1$. Fix an index $i$ and drop subscripts: set
$n=n_i$, $k=k_i$, $V=V_i$, and so on.

\paragraph{Construction of $H$ (\cite[\S 3]{FPS}).}
Choose a real parameter $\delta>1$ and set $d:=\delta k$. Let
$U(d)\subseteq V$ be the set of vertices of degree at least~$d$, capped
at $k-\varphi(k)$, where $\varphi(k):=k^{0.9}$ is the slack term of
\cite[Footnote 1, p.~6]{FPS}. Add to $U(d)$ a uniformly random set of
$k-\ell$ vertices from $V\setminus U(d)$, where $\ell=|U(d)|$, to form
$U$. After a deterministic edge-swap pre-processing
(\cite[p.~7--8]{FPS}) and an additional uniform-random completion of
the maps $f_u\colon U\setminus(N(u)\cup\{u\})\to N(u)\cap W$, one
obtains the multigraph $H$ on $W:=V\setminus U$. By
Vizing--Gupta~\cite[Lem.~2.2(ii)]{FPS},
\begin{equation}\label{eq:VG}
  \chip(H)\;\le\;\Delta(H)+\mu(H).
\end{equation}
The bound~\eqref{eq:fpsbound} reduces (\cite[\S 3]{FPS}) to two
high-probability estimates:
\begin{align}
  &\textbf{Proposition 3.3.}\quad
  \Delta(H)\le \bigl(F(\delta)+o(1)\bigr)\,k,\label{eq:prop33}\\[2pt]
  &\textbf{Proposition 3.4.}\quad
  \mu(H)\;=\;o(k),\label{eq:prop34}
\end{align}
where $F(\delta)$ is the upper-bound function whose value $F(9/8)=9/16$
is established by~\cite[Claim 3.7]{FPS}.

\paragraph{The optimisation (\cite[Claim 3.7]{FPS}).}
The function $F(\delta)$ is the supremum of the asymptotic
expression
\begin{equation}\label{eq:obj}
   f(\alpha,\beta,\gamma;\delta)
   \;=\;
   \gamma - \alpha\,\frac{\delta-1}{\delta-\gamma},
\end{equation}
over the feasible region
\begin{equation}\label{eq:constraints}
   0\le\alpha\le\beta\le 1,
   \qquad
   0\le\gamma\le\alpha+(\delta-\alpha)\,\frac{1-\beta}{2-\beta}.
\end{equation}
Here $\alpha,\beta,\gamma$ are the normalisations
$\alpha=\ell_w/k$, $\beta=\ell/k$, $\gamma=|U_w|/k$,
where $\ell_w$ is the number of $U(d)$-neighbours of a target vertex
$w\in W$ and $U_w:=N(w)\cap U$.

Following~\cite[Proof of Claim 3.7]{FPS}, the supremum splits into
three cases:
\begin{description}
\item[Case 1.] $\alpha=0$.
\item[Case 2a.] The upper constraint on $\gamma$
  in~\eqref{eq:constraints} is tight, and the stationary point of $f$
  in $\alpha$ is interior to $[0,\beta]$.
\item[Case 2b.] The upper constraint on $\gamma$ is tight, but the
  stationary point falls outside $[0,\beta]$, forcing $\alpha=\beta$.
\end{description}
Reparametrising $\eta:=1-(1-\beta)/(2-\beta)\in[\tfrac12,1]$
(equivalently $\beta=2-1/\eta$), Case~2 has
$\gamma=\eta\alpha+(1-\eta)\delta$.

\section{The optimisation with \texorpdfstring{$\delta$}{delta} free}
\label{sec:optimisation}

We compute the maximum of $f$ in each case as an explicit function of
$\delta$. All three formulas reduce to simple expressions and lead to
the cubic~\eqref{eq:cubic}.

\begin{lemma}[Case 1]\label{lem:case1}
For every $\delta\ge 1$, the maximum of $f$ on the feasible set
restricted to $\alpha=0$ is
\[
  f_1(\delta) \;=\; \frac{\delta}{2},
\]
attained at $\beta=0,\gamma=\delta/2$.
\end{lemma}

\begin{proof}
With $\alpha=0$, \eqref{eq:obj} gives $f=\gamma$, and the upper bound
on $\gamma$ in~\eqref{eq:constraints} simplifies to
$\gamma\le\delta(1-\beta)/(2-\beta)$. This is a decreasing function of
$\beta\in[0,1]$ with derivative $-1/(2-\beta)^2<0$, hence its maximum
over $\beta\in[0,1]$ is at $\beta=0$, with value $\delta/2$.
\end{proof}

\begin{lemma}[Case 2 stationary point and 2a/2b boundary]\label{lem:stationary}
For every $\delta>1$ and $\eta\in[\tfrac12,1]$, the unique stationary
point in $\alpha$ of $f$ on the Case-2 line
$\gamma=\eta\alpha+(1-\eta)\delta$ is
\begin{equation}\label{eq:alphastar}
   \alpha^{\star}(\eta,\delta)
   \;=\;
   \delta - \frac{\sqrt{\delta(\delta-1)}}{\eta}.
\end{equation}
This stationary point equals $\beta(\eta)=2-1/\eta$ exactly when
\begin{equation}\label{eq:etaboundary}
   \eta=\eta_{b}(\delta)
   \;:=\;\frac{\sqrt{\delta(\delta-1)}-1}{\delta-2}.
\end{equation}
The Case 2a regime is $\eta\in[\eta_b(\delta),1]$ and Case 2b is
$\eta\in[\tfrac12,\eta_b(\delta))$. At $\delta=9/8$ this recovers
\cite[(p.~10)]{FPS}: $\alpha^{\star}=(9-3/\eta)/8$ and
$\eta_b=5/7$.
\end{lemma}

\begin{proof}
Direct differentiation of
$f(\alpha,\cdot,\eta\alpha+(1-\eta)\delta;\delta)$ in $\alpha$ and
setting the result to zero yields a quadratic in $\alpha$ with the two
roots $\delta\pm\sqrt{\delta(\delta-1)}/\eta$. Only
$\alpha^{\star}=\delta-\sqrt{\delta(\delta-1)}/\eta$ lies in $[0,\beta]$
(at $\delta=9/8$ and $\eta=3/4$ it gives $\alpha^{\star}=5/8$, while the
other root gives $\alpha=5/3>1$). Equating
$\alpha^{\star}(\eta,\delta)=2-1/\eta$ and solving for $\eta$
gives~\eqref{eq:etaboundary}. The substitutions at $\delta=9/8$ are
elementary.
\end{proof}

\begin{lemma}[Case 2a]\label{lem:case2a}
For every $\delta>1$, $f$ restricted to Case 2a is strictly decreasing
in $\eta$ on $[\eta_b(\delta),1]$; its maximum is therefore attained at
$\eta=\eta_b(\delta)$. Setting
$\alpha=\alpha^{\star}(\eta_b(\delta),\delta)$ and simplifying,
\begin{equation}\label{eq:case2amax-clean}
   f_{2a}^{\max}(\delta)
   \;=\;
   \frac{-\delta^{5/2}-\delta^{3/2}+2\sqrt{\delta}
         + (\delta^2+2\delta-2)\sqrt{\delta-1}}
        {-\delta^{3/2}+\sqrt{\delta}+\sqrt{\delta-1}}.
\end{equation}
At $\delta=9/8$, $f_{2a}^{\max}(9/8)=11/20$, matching~\cite[p.~10]{FPS}.
\end{lemma}

\begin{remark}\label{rem:case2a-form}
The closed form~\eqref{eq:case2amax-clean} is what one obtains by
substituting the stationary point~\eqref{eq:alphastar} into the Case-2
expression $\gamma=\eta\alpha+(1-\eta)\delta$ inside~\eqref{eq:obj},
then setting $\eta=\eta_b(\delta)$. The function is smooth on
$\delta>1$, tends to $1$ as $\delta\to1^+$, has a unique minimum near
$\delta\approx1.20$ with value $\approx 0.534$, and tends to $+\infty$
as $\delta\to+\infty$.
\end{remark}

\begin{lemma}[Case 2b]\label{lem:case2b}
For every $\delta>1$, the maximum of $f$ over Case 2b is attained at
$\eta=\tfrac12$ (so $\alpha=\beta=0$) and equals
\[
   f_{2b}^{\max}(\delta) \;=\; \frac{\delta}{2}.
\]
\end{lemma}

\begin{proof}
At $\eta=\tfrac12$, $\alpha=\beta=2-1/\eta=0$; from
$\gamma=\eta\alpha+(1-\eta)\delta=\delta/2$ and
$f=\gamma-\alpha(\delta-1)/(\delta-\gamma)$, $f=\delta/2$ since
$\alpha=0$.

The verification that $\eta=1/2$ is the maximum (i.e. that
$f$ is monotone decreasing in $\eta$ on $[\tfrac12,\eta_b(\delta))$
under the substitution $\alpha=\beta=2-1/\eta$) is a one-variable
calculus check; at $\delta=9/8$ it matches~\cite[p.~10]{FPS} verbatim.
\end{proof}

\begin{proposition}\label{prop:F-form}
For every $\delta>1$,
\[
   F(\delta) \;=\; \max\bigl(\,f_1(\delta),\,f_{2a}^{\max}(\delta),\,
                              f_{2b}^{\max}(\delta)\bigr)
            \;=\; \max\!\bigl(\tfrac{\delta}{2},\,f_{2a}^{\max}(\delta)\bigr).
\]
At $\delta=9/8$, $f_{2a}^{\max}(9/8)=11/20<\delta/2=9/16$, so $F(9/8)=9/16$,
matching~\cite[Claim 3.7]{FPS}.
\end{proposition}

\begin{proof}
Lemmas~\ref{lem:case1},~\ref{lem:case2b} give $f_1=f_{2b}^{\max}=\delta/2$.
\end{proof}

\begin{lemma}\label{lem:cubic}
Let $\delta_1$ be the unique real root in $(1,5/4)$ of the irreducible
cubic
\(
   \delta^3+3\delta^2-\delta-4=0.
\)
Then the equation $\delta/2=f_{2a}^{\max}(\delta)$ has three real
solutions in $(1,2)$: $\delta=\delta_1$, $\delta=4/3$, and an extraneous
solution at $\delta=1$ arising from the rationalisation; the cubic
$\delta^3+3\delta^2-\delta-4$ has its other two real roots in
$\RR_{<0}$, hence not in $(1,2)$.
\end{lemma}

\begin{proof}
Multiplying the equation $\delta/2=f_{2a}^{\max}(\delta)$ through and
substituting $s:=\sqrt\delta$, $t:=\sqrt{\delta-1}$ collects every
non-rational term into a single product $s^j t$; rearranging gives an
identity of the form $A(s)=B(s)\,t$, and squaring (using $t^2=s^2-1$)
yields a polynomial identity in $s$ alone. After substituting $s^2=\delta$,
the resulting degree-5 polynomial factorises (over $\QQ$) as
\begin{equation}\label{eq:factor}
   -\tfrac{1}{4}\,(\delta-1)\,(3\delta-4)\,(\delta^3+3\delta^2-\delta-4).
\end{equation}
The linear factors give $\delta=1$ (extraneous, an artefact of the
squaring) and $\delta=4/3$. The cubic factor has discriminant
\[
   \mathrm{disc}(\delta^3+3\delta^2-\delta-4) \;=\; 229 \;>\;0,
\]
so it has three distinct real roots; numerical evaluation gives
$\delta_1\approx1.1149075414767557$,
$\delta_2\approx-1.2541016883650524$, and
$\delta_3\approx-2.8608058531117035$. Irreducibility follows from the
rational-root test: the only rational candidates $\pm1,\pm2,\pm4$ are
checked not to be roots.
\end{proof}

\begin{theorem}\label{thm:Foptimum}
$F$ attains its global minimum on $(1,\infty)$ at $\delta=\delta_1$,
with value
\[
   F^{\star} \;:=\; F(\delta_1) \;=\; \frac{\delta_1}{2}.
\]
\end{theorem}

\begin{proof}
By~\Cref{prop:F-form}, $F(\delta)=\max(\delta/2,f_{2a}^{\max}(\delta))$.
The function $\delta/2$ is strictly increasing on $\delta>1$. The
function $f_{2a}^{\max}(\delta)$ is U-shaped on $\delta>1$: it tends to
$1$ as $\delta\to1^+$, is strictly decreasing on
$(1,\delta_{\mathrm{opt}})$ for some $\delta_{\mathrm{opt}}>9/8$, and
strictly increasing on $(\delta_{\mathrm{opt}},\infty)$ (this last
statement is a one-variable calculus exercise on
$f_{2a}^{\max}(\delta)$ but is most economically verified numerically,
see~\Cref{sec:numerics}).

For $\delta\in(1,\delta_1)$, $f_{2a}^{\max}(\delta)>\delta/2$, so
$F(\delta)=f_{2a}^{\max}(\delta)$ is decreasing. For
$\delta>\delta_1$, $f_{2a}^{\max}(\delta)\le\delta/2$ until the next
intersection point at $\delta=4/3$
(\Cref{lem:cubic}); on this whole interval
$F(\delta)=\delta/2$ is increasing. Hence the global minimum of $F$ on
$(1,\infty)$ is at the first intersection point
$\delta=\delta_1$, with $F^{\star}=\delta_1/2$.
\end{proof}

\begin{proof}[Proof of~\Cref{thm:main}]
By~\Cref{thm:Foptimum}, $F^{\star}=\delta_1/2<9/16$. Substituting
$\delta=2x$ in the cubic of~\Cref{lem:cubic} gives
\begin{equation}\label{eq:cubic-in-F}
   (2x)^3+3(2x)^2-(2x)-4
   \;=\;
   8x^3+12x^2-2x-4
   \;=\;
   2\,(4x^3+6x^2-x-2),
\end{equation}
so $F^{\star}=\delta_1/2$ is the unique real root of $4x^3+6x^2-x-2=0$
in $(\tfrac12,\tfrac58)$. To put $\delta_1$ in trigonometric form,
substitute $\delta=y-1$ in the cubic of~\Cref{lem:cubic}; this yields
the depressed cubic
\begin{equation}\label{eq:depressed}
   y^3 - 4y - 1 \;=\; 0.
\end{equation}
With $p=-4$, $q=-1$, the discriminant is $-4p^3-27q^2 = 256-27=229>0$,
so all three roots are real; Cardano's trigonometric formula gives
\[
   y_k \;=\; \frac{4}{\sqrt3}\,
            \cos\!\left(\frac13\arccos\!\frac{3\sqrt3}{16}-\frac{2\pi k}{3}\right),
   \qquad k\in\{0,1,2\}.
\]
The principal value $y_0=(4/\sqrt3)\cos((1/3)\arccos(3\sqrt3/16))$
satisfies $y_0\approx2.115$, hence
$\delta_1=y_0-1\approx1.115\in(1,5/4)$ as desired, and
$F^{\star}=\delta_1/2$ takes the closed form stated
in~\Cref{thm:main}.

The remainder of the proof of~\eqref{eq:fpsbound} -- specifically
Propositions~3.3 and~3.4 of~\cite{FPS} and the supporting Claims~3.5
and~3.6 -- is robust to the replacement $\delta=9/8\rightsquigarrow\delta_1$,
as verified in~\Cref{sec:robustness}. Hence Lemma~2.3 of~\cite{FPS}
holds with $\cFPS$ replaced by~$F^{\star}$.
\end{proof}

\section{Robustness of \texorpdfstring{\cite{FPS}}{[FPS]}, Propositions 3.3 and 3.4}
\label{sec:robustness}

The proof of~\eqref{eq:fpsbound} uses, beyond Claim~3.7 itself, the
following auxiliary lemmas of~\cite{FPS}:
\begin{itemize}
\item Claim~3.5 (\cite[p.~8]{FPS}): upper bound on $|U_w|$ given
  $\ell_w,\ell$.
\item Claim~3.6 (\cite[p.~8]{FPS}): high-probability lower bound on
  $\delta\cdot k - |U_w|$.
\item Proposition~3.3 (\cite[p.~8]{FPS}): $\Delta(H)\le(F(\delta)+o(1))k$.
\item Proposition~3.4 (\cite[p.~11]{FPS}): $\mu(H)=o(k)$.
\end{itemize}

\Cref{thm:Foptimum} is essentially a restatement of Proposition~3.3
with $\delta$ free, so the latter holds at $\delta=\delta_1$ by
construction. We check the remaining three.

\paragraph{Claim 3.5.}
Reads: $|U_w|\le \ell_w+(d-\ell_w)(k-\ell)/(2k-\ell-1)+o(k)$. Its
proof is a counting argument that uses no numerical specialisation of
$d/k$ beyond $d\ge\ell+1$, which is satisfied for any $\delta\in(1,5/4)$.

\paragraph{Claim 3.6.}
States that the random part of $U$ contributes at most a certain
expected number of additional $U$-neighbours to $w$, with an error
$o(k)$. Reading~\cite[p.~9]{FPS}: ``Set $4d<5k$.'' This is the only
place in the proof of Claim~3.6 where $\delta=d/k$ is explicitly
constrained, requiring $\delta<5/4$. Since
$\delta_1\approx 1.115<5/4=1.25$, Claim~3.6 holds at $\delta=\delta_1$
with the same proof.

\paragraph{Proposition 3.4.}
The proof (\cite[p.~11]{FPS}) uses $d=\delta k$ in the single
quantitative bound
\[
   \#\bigl\{w\in W: w \text{ has }\ge\varphi(k)\text{ neighbours in }U\bigr\}
   \;\le\; \frac{kd}{\varphi(k)}
   \;\le\; 2\,k^{1.1},
\]
i.e.\ the assertion $\delta\le 2$. This is trivially satisfied at
$\delta_1\approx1.115$. The rest of the Proposition 3.4 proof uses
$\delta$ only through the cap $|U(d)|\le k-\varphi(k)$ in~\cite[Footnote
1, p.~6]{FPS}, which is a structural choice independent of the
numerical value of $\delta$.

The aggregate admissible range for $\delta$ across all of these inputs
is $\delta\in(1,5/4)$, and $\delta_1=1.114907\ldots$ lies in the
interior of this interval.

\section{Numerical verification and reproducibility}\label{sec:numerics}

We verified the calculations of~\Cref{sec:optimisation} symbolically
using SymPy~1.14. The script reproduces $F(9/8)=9/16$,
$f_{2a}^{\max}(9/8)=11/20$, $\alpha^{\star}(\eta,9/8)=(9-3/\eta)/8$,
$\eta_b(9/8)=5/7$ exactly, and computes the rationalised polynomial
in~\eqref{eq:factor}, the discriminant $229$ of the cubic, and the
roots of the cubic and of $4x^3+6x^2-x-2$ to 30 significant figures.
A second script verifies that the FPS-printed Case~2b objective
$(7\eta+1)/8-(2+1/\eta)/(8-7\eta)$ (\cite[p.~10]{FPS}) appears to
contain a sign typo: replacing the $+$ by $-$ in the numerator of the
second term gives exactly the value $9/16$ at $\eta=1/2$, while the
printed form does not.

\begin{remark}[Sign typo in {\cite[p.~10]{FPS}}]\label{rem:sign}
The change $(2+1/\eta)\rightsquigarrow(2-1/\eta)$ in the second term
of the Case 2b objective is consistent with the surrounding algebra
($\alpha=2-1/\eta$, $\delta=9/8$, hence $\alpha(\delta-1)=(2-1/\eta)/8$)
and gives the claimed value $9/16$ at $\eta=1/2$. The printed sign
gives a negative value at $\eta=1/2$, contradicting~\cite[Claim 3.7]{FPS}.
The conclusion of Claim 3.7 is unaffected.
\end{remark}

\section{Downstream and limitations}\label{sec:remarks}

\begin{remark}[Downstream impact on \cite{FPS} Theorem 1.2]\label{rem:downstream}
Theorem~1.2 of~\cite{FPS} states that Albertson's
conjecture~\eqref{eq:albertson} holds for every $k$-chromatic graph on
at most $1.4(k-1)$ vertices, with an asymptotic bound of
$(1.64-o(1))\,k$. The asymptotic coefficient $1.64$ is derived
from $\cFPS=9/16$ via the inequality
$\chip(H_i)\le k-k_i$. Replacing $\cFPS$ by $F^{\star}$
(\Cref{thm:main}) replaces $1.64$ by a slightly larger value
$g(F^{\star})$, where $g$ is the explicit function of $\cFPS$ given
in~\cite[\S 2]{FPS}; the magnitude of the change is on the order of
$10^{-2}$, i.e.\ the asymptotic Albertson bound moves from
$(1.64-o(1))\,k$ to roughly $(1.65-o(1))\,k$. We do not redo the
bookkeeping here. The improvement leaves the finite Cranston
residuals $(t,|G|)\in\{(25,48),(26,50),(26,51)\}$~\cite{Cranston}
untouched.
\end{remark}

\begin{remark}[Sharpness of $F^{\star}$]\label{rem:sharp}
$F^{\star}$ is sharp \emph{within} the three-case framework
of~\cite[Claim 3.7]{FPS}, the Vizing--Gupta bound~\eqref{eq:VG}, and
the semi-random construction of~\cite[\S 3]{FPS}. Several routes might
yield further improvement:
\begin{enumerate}
\item Replacing Vizing--Gupta in~\eqref{eq:VG} by a stronger
  chromatic-index estimate. For instance, the Goldberg--Seymour
  theorem (\cite{ChenJingZang} for the proof) bounds
  $\chip(H)$ by $\max(\Delta(H), \lceil\Gamma(H)\rceil)$, where
  $\Gamma(H)$ is the fractional chromatic index; replacing
  $\Delta+\mu$ by $\max(\Delta,\lceil\Gamma\rceil)$ removes the
  $\mu$ contribution entirely if $\Gamma\le\Delta$, which is conjectural
  but plausible for the semi-random multigraphs constructed in~\cite{FPS}.
\item Refining the case split in~\cite[Claim 3.7]{FPS}. The Case-2b
  binding inequality reduces to the elementary estimate
  $|U_w|\le d/2$ when the high-degree set $U(d)$ is empty; this is
  tight under uniform random selection and cannot be improved without
  changing the random selection itself.
\item Changing the semi-random construction in~\cite[\S 3]{FPS} so
  that the threshold $d=\delta k$ is replaced by a non-constant or
  adaptive choice.
\end{enumerate}
Each of these is a separate research direction.
\end{remark}

\section*{Acknowledgements}
This note is a small follow-up to the work of Fox, Pach and Suk~\cite{FPS}
and would not exist without their construction. Computer-algebra
verification was performed in SymPy~1.14.

\bibliographystyle{plain}
\begin{thebibliography}{99}

\bibitem{Albertson}
M.~O. Albertson.
Posted at
\url{http://www.openproblemgarden.org/op/crossing_numbers_and_coloring},
2007/2009.

\bibitem{Ackerman}
E.~Ackerman.
On topological graphs with at most four crossings per edge.
\textit{Computational Geometry: Theory and Applications}, 2019.
\texttt{arXiv:1509.01932}.

\bibitem{ACF}
M.~O. Albertson, D.~W. Cranston, and J.~Fox.
Crossings, colorings, and cliques.
\textit{Electronic Journal of Combinatorics}, 16(1):R45, 2009.
\texttt{arXiv:1006.3783}.

\bibitem{BaratToth}
J.~Bar\'at and G.~T\'oth.
Towards the Albertson conjecture.
Preprint, 2009.
\texttt{arXiv:0909.0413}.

\bibitem{ChenJingZang}
G.~Chen, G.~Jing, and W.~Zang.
Proof of the Goldberg--Seymour conjecture on edge-colorings of multigraphs.
Preprint, 2019.
\texttt{arXiv:1901.10316}.

\bibitem{Cranston}
D.~W. Cranston.
Progress on Albertson's conjecture.
Preprint, 2025.
\texttt{arXiv:2512.08020}.

\bibitem{FPS}
J.~Fox, J.~Pach, and A.~Suk.
Immersions and Albertson's conjecture.
In \textit{41st International Symposium on Computational Geometry (SoCG 2025)}.
LIPIcs vol.~332, Schloss Dagstuhl, 2025.
\texttt{arXiv:2510.05893}.

\end{thebibliography}

\end{document}
